\documentclass[11pt]{article}

% basic packages
\usepackage[margin=1in]{geometry}
\usepackage[pdftex]{graphicx}
\usepackage{amsmath,amssymb,amsthm}
\usepackage{william}

% page formatting
\usepackage{fancyhdr}
\pagestyle{fancy}

\renewcommand{\sectionmark}[1]{\markright{\textsf{\arabic{section}. #1}}}
\renewcommand{\subsectionmark}[1]{}
\lhead{\textbf{\thepage} \ \ \nouppercase{\rightmark}}
\chead{}
\rhead{}
\lfoot{}
\cfoot{}
\rfoot{}
\setlength{\headheight}{14pt}

\linespread{1.03}
\setlength{\parindent}{0.2in}
\setcounter{secnumdepth}{2}
\setcounter{tocdepth}{2}

\begin{document}

\thispagestyle{empty}
\bigskip \
\vspace{0.1cm}

\begin{center}
{\fontsize{22}{22} \selectfont Lecture Notes on}
\vskip 16pt
{\fontsize{36}{36} \selectfont \bf \sffamily Condensed Matter}
\vskip 24pt
{\fontsize{18}{18} \selectfont \rmfamily \href{https://will-lancer.github.io}{Will Lancer}}
\vskip 6pt
{\fontsize{14}{14} \selectfont \ttfamily will.m.lancer@gmail.com}
\vskip 24pt
\end{center}

{\parindent0pt \baselineskip=15.5pt}
\noin
cond-mat

\newpage
\microtoc
\newpage

\section{Lattices}

\subsection{Lattice structures}

The atoms in most solids arrange themselves in periodically repeating
structures called \emph{lattices}. Formally we'll say that an $m$-dimensional lattice
$\Lambda$ is a $\ints$-module generated by some linearly independent subset of vectors 
$\mathbf{a}_1, \ldots, \mathbf{a}_m \in \reals^n$. We call these vectors \emph{primitive
lattice vectors}. The primary example to keep in mind is the \emph{square lattice},
which is the lattice generated by the standard basis vectors, $\mathbf{a}_i = \bfe_i$.

Given a lattice $\Lambda \subset \reals^n$, we can mod out by the lattice structure to implement 
the physical equivalence
\begin{align*}
    \bfx \sim \bfx + \sum_i n_i \mathbf{a}_i.
\end{align*}
This gives us the ``fundamental domain'' of our space, i.e. the space of actually
novel motions in it. We'll denote this as $X = \reals^n/\Lambda$. This will be useful
later when we talk about functions on the lattice.

We can consider portions of our space at a time---we call these subsets of 
$\reals^n$ \emph{cells}. There are three kinds of cells that are disproportionately 
useful to consider. First, notice that if we consider $X$ as a subset of $\reals^n$, 
by translating it around with $\Lambda$, we fill the entire space $\reals^n$. Said
differently, the orbit of $X$, $\bigo_X = \reals^n$. We will call
cells of this type \emph{unit} cells. We will call the unit cell with the fewest
generators a \emph{primitive} cell. Lastly, there is a special kind of unit
cell that naturally manifests the symmetries of the lattice $\Lambda$
called a \emph{Wigner-Seitz} cell. It's defined as
\begin{align*}
    \Gamma = \{ \bfx \in \reals^n \mid |\bfx| < |\bfx - \mathbf{A}| \text{ for all } \mathbf{A} \in \Lambda \}.
\end{align*}
That is: define some arbitrary point in your lattice as the origin
and consider the closest elements to it such that translating the cell
around fills the entire space---this is the Wigner-Seitz cell.
We can see that this cell naturally manifests the symmetries of the
lattice by considering the group action of $\Lambda \circlearrowright \Gamma$.
Letting $g \in \Lambda$ act on the inequality above, we have
\begin{align*}
    |g\bfx| < |g\bfx - g \mathbf{A}| \text{ for all } \mathbf{A} \in \Lambda. 
\end{align*}
But, $g\mathbf{A}$ runs over the elements of $\Lambda$ when $\mathbf{A}$
does, and thus if $\bfx \in \Gamma$, then
\begin{align*}
    |g\bfx| < |g\bfx - \mathbf{A}| \text{ for all } \mathbf{A} \in \Lambda. 
\end{align*}
That is: $g\bfx \in \Gamma$ as well, so $g \Gamma = \Gamma$,
i.e. the Wigner-Seitz cell only rearranges itself under the
lattice symmetries.

There are a few natural constructions we can now take stock of. 
The first is to note that we can easily define the lattice dual in the standard way,
\begin{align*}
    \Lambda^\vee = \Hom(\Lambda, \reals).
\end{align*}
We can define a basis for $\Lambda^\vee$ in natural way,
i.e. by considering the set of functionals $\varphi \in \Lambda^\vee$
dual to the primitive lattice vectors,
\begin{align*}
    \varphi^i(\mathbf{a}_j) = \delta^i_j.
\end{align*}
Furthermore, since $\reals^n$ has an inner product on it, we
can associate to each functional $\varphi^i$ a unique vector $\mathbf{g}_i \in \Lambda$
such that for any $\mathbf{A} \in \Lambda$,
\begin{align*}
    \varphi^i(\mathbf{A}) = \mathbf{g}_i \cdot \mathbf{A}.
\end{align*}
This is a natural consequence of the isomorphism $\Lambda \cong \Lambda^\vee$.
For our purposes, it will be more useful to rescale the reciprocal
basis vectors by $2\pi$; that is, instead choose a basis of $\Lambda^\vee$
such that
\begin{align*}
    \mathbf{g}_i \cdot \mathbf{a}_j = 2\pi \delta_{ij}.
\end{align*}
$\Lambda^\vee$ itself can be thought of as a lattice due to the isomorphism
$\Lambda^\vee \cong \Lambda$, and so all of the previous discussion about
cells goes through the same. We call the Wigner-Seitz cell for $\Lambda^\vee$
the \emph{Brillouin zone}, $\mathcal{B}$, and it's one of the most important constructions
in condensed matter. 

For example, if $\Lambda = \langle a \rangle$, then the Wigner-Seitz cell for
the lattice is $\Gamma = [-a/2, a/2)$, and the Brillouin zone is $\mathcal{B} = [-\pi/a, \pi/a)$.
This is easily seen by noting that $\Lambda^\vee = \langle g = 2\pi/a \rangle$,
as $g a = 2\pi$, and so $\mathcal{B} = [-g/2, g/2)$---an line of diameter $g$
centered at the origin.



\subsection{Constructions on the lattice}

We now define functions $\psi(\bfx)$ on $\reals^n$ with an embedded lattice
$\Lambda \hookrightarrow \reals^n$. It's intuitively clear that
if we describe $\psi(\bfx)$ on $X = \reals^n/\Lambda$, then by
translational symmetry we've described it on all of $\reals^n$.

However, all of the above discussion has implicitly been assuming
a free particle $\psi(\bfx)$---\emph{Bloch's theorem} gives us insight
into what changes when we add a potential $V(\bfx)$ that respects the lattice
structure. Let $V(\bfx) = V(\bfx + \mathbf{A})$, where $\mathbf{A} \in \Lambda$;
that is, $V$ is a periodic potential on the lattice. Then Bloch's theorem
states that we can label energy eigenstates purely by their crystal momentum $\mathbf{k}$:
\begin{align*}
    \boxed{\psi_\bfk(\bfx) = e^{i\bfk \cdot \bfx} u_\bfk(\bfx),}
\end{align*}
where $u_\bfk(\bfx)$ is a function periodic in $\Lambda$. We have
been implicitly working in the \emph{extended zone scheme} for $\bfk$.
This is where we do not implement the equivalence $\bfk \sim \bfk + \mathbf{G}$,
with $G \in \Lambda^\vee$, and so $\bfk$ ranges through all of $\reals^n$.
We may also explicitly implement this equivalence and work in the
\emph{reduced zone scheme}, where we let $\bfk \in \mathcal{B}$ instead,
and then add an extra index to $\psi$ to resolve degeneracies:
\begin{align*}
    \psi_{\bfk, n}(\bfx) = e^{i\bfk \cdot \bfx} u_{\bfk, n}(\bfx).
\end{align*}
The way to prove this is as follows: consider a representation of $\Lambda$,
$\rho_\Lambda \colon \Lambda \to GL(\mathcal{H})$ by $\mathbf{A} \mapsto T_\mathbf{A}$,
where $\mathcal{H}$ is our Hilbert space and $\mathbf{A} \in \Lambda$.
These are \emph{translation operators} $T_\mathbf{A}$ acting on $\psi$ as
\begin{align*}
    T_{\mathbf{A}} \psi(\bfx) = \psi(\bfx + \mathbf{A}).
\end{align*}
It is trivial to verify that the $T$'s form a representation of $\Lambda$.
Note that since $\Lambda$ is abelian, $T_{\mathbf{A}} = e^{i\theta_\mathbf{A}}$
for some $\theta_\mathbf{A} \in [0, 2\pi)$.

Since the wave functions are symmetric under translations by $\Lambda$,
the translation operators commute with the Hamiltonian, $[T_\mathbf{A}, H] = 0$.
Thus, $H$ and $T_\mathbf{A}$ are simultaneously diagonalizeable. We now have
to consider what the natural quantum number is to label our wave-functions
with. Recall that, since translation is a commutative operation, the translation
operators must act like a phase. Furthermore, since $\psi(\bfx) = \psi(\bfx + \mathbf{A})$
for any $\mathbf{A} \in \Lambda$, we must have
\begin{align*}
    \psi(\bfx) = \psi(\bfx + \mathbf{A}) = T_{\mathbf{A}} \psi(\bfx) = e^{i\theta_{\mathbf{A}}} \psi(\bfx)
    \implies \theta_\mathbf{A} = 2\pi n.
\end{align*}
Thus, we can write $\theta_{\mathbf{A}} = \bfk \cdot \mathbf{A}$, where $\bfk \in \Lambda^\vee$,
as $\bfk = \sum_i a_i \mathbf{g}_i$, and $\mathbf{g}_i \cdot \mathbf{a}_j = 2\pi \delta_{ij}$.
So, we can label our states by the quantum number $\bfk \in \mathcal{B}$.
Bloch's theorem follows when one considers the definition of $u_\bfk(\bfx) = e^{-i\bfk \cdot \bfx} \psi_\bfk(\bfx)$
and implements the periodicity of the lattice structure. 












\end{document}